% registry: ZHANGYAO
\methodhead{ZHANGYAO}{Zhang and Yao (2011): nonlinear critical-plane fatigue-limit criterion}\label{sheet:ZHANGYAO}

\block{Source} C.-C.~Zhang, W.-X.~Yao, \emph{Fatigue Fract.\ Engng Mater.\ Struct.} 34 (2011)
337--344 \cite{zhang2011}, Eqs.~(3), (6)--(10), (A1); Tables~2--8.

\block{Idea and scope} As in McDiarmid's criterion, the critical plane is the plane of maximum shear
stress amplitude $C_a$. Unlike McDiarmid's linear sum, the shear amplitude and the normal stress $N^*$
on that plane are combined nonlinearly, so that both fully reversed limits are met exactly; this
extends the criterion to $1.25<\fm/\tm<2$. The mean normal stress on the critical plane enters by a
Goodman correction of the normal stress amplitude. Proposed for high-cycle fatigue limits of ductile
and hard metals under proportional and out-of-phase tension (or bending) and torsion, with mean
stresses.

\block{Inputs}
\begin{center}\small
\begin{tabular}{@{}llp{100mm}@{}}
\toprule
Symbol & Unit & Meaning\\\midrule
$(\varphi_c,\theta_c)$ & rad & orientation of the critical plane\\
$C_a$ & MPa & shear stress amplitude on the plane: major semi-axis of the shear ellipse (Eq.~A1)\\
$N_a$, $N_m$ & MPa & normal stress amplitude and mean on the plane\\
$N^*$ & MPa & Goodman-corrected normal stress amplitude, Eq.~\eqref{eq:zy-nstar}\\
$\sigma_u=S_u$ & MPa & ultimate tensile strength (workbook \texttt{Su})\\
$\fm$, $\tm$ & MPa & fully reversed bending (or tension) and torsion fatigue limits\\
LHS & MPa & left-hand side of Eq.~\eqref{eq:zy}\\
$I$ & \% & error index, Eq.~(10)\\
$\mathrm{DP}$ & -- & repository index, LHS$/\tm$, Eq.~\eqref{eq:zy-dp}\\
\bottomrule
\end{tabular}
\end{center}

\block{Formulation} On the plane $(\varphi_c,\theta_c)$ of maximum $C_a$ (Eqs.~6--7),
\begin{gather}
\sqrt{C_a^2 - C_a N^* + \frac{\tm^2}{\fm^2}\bigl(\fm + 2C_a\bigr)N^*} \le \tm ,
\label{eq:zy}\\
N^* = \frac{N_a}{1-N_m/\sigma_u} .
\label{eq:zy-nstar}
\end{gather}
Fully reversed torsion gives $C_a=\tm$, $N^*=0$; fully reversed bending gives $C_a=N^*=\fm/2$ on the
$45^\circ$ plane; both make the left-hand side equal to $\tm$. The condition that a compressive mean be
beneficial (Eqs.~8--9) leads to $\fm/\tm<2$; the lower bound 1.25 is taken from Papadopoulos. The error
index is $I=100\,(\mathrm{LHS}-\tm)/\tm$ (Eq.~10). The repository evaluates on every plane
\begin{equation}
\mathrm{DP} = \frac{1}{\tm}\sqrt{C_a^2 - C_a N^* + \frac{\tm^2}{\fm^2}(\fm+2C_a)N^*},
\qquad I = 100\,(\mathrm{DP}-1).
\label{eq:zy-dp}
\end{equation}
\emph{Critical plane}: \texttt{ZHANGYAO} takes the plane of largest $C_a$; \texttt{ZHANGYAO\_ext} the
largest DP over all planes. \emph{Ties} of the maximum $C_a$ go to the larger DP. \emph{Shear
amplitude}: half the longest chord of the shear path (\texttt{longest\_chord}, MCC), the major
semi-axis of Eq.~(A1), by default; the global selection of \texttt{docs/shear-amplitude.md} can switch
it to \texttt{rss} (MCE).

\block{Constants and calibration} Paper and repository: nothing is fitted. Eq.~\eqref{eq:zy} is built so
that $\fm$ and $\tm$ lie on it, and $\sigma_u$ enters only the mean-stress correction. The repository
reads workbook \texttt{Sigm1}, \texttt{Taum1}, \texttt{Su}; only the master-curve constants are fitted,
on the uniaxial $R=-1$ design tests.

\block{Life} The paper gives none; finite lives use the master curve on the index
(Section~\ref{sec:mastercurve}), a repository extension. On fatigue-limit datasets $\mathrm{DP}-1$ is
the paper's $I/100$.

\block{Repository choices}
\begin{itemize}
\item $C_a$ by the longest-chord rule; the tables are not reproduced with \texttt{rss}.
\item Ties of the maximum $C_a$ towards the larger DP; the paper gives no rule, and this criterion is
  sensitive to it (see Verification).
\item $\sigma_u$ is needed only when a plane carries a mean normal stress; without it such a case is
  skipped (\texttt{MissingParamsError}). $N_m\ge\sigma_u$, or a negative radicand in
  Eq.~\eqref{eq:zy}, is out of scope (\texttt{LoadingScopeError}).
\item A warning is issued outside $1.25<\fm/\tm<2$.
\item The fatigue-limit index is used as the damage parameter of the shared master curve.
\end{itemize}

\block{Worked example} Hard steel, Table~2: $\fm=313.9$, $\tm=196.2$, $\sigma_u=680\MPa$; 42CrMo4:
$\fm=398$, $\tm=260$, $\sigma_u=1025\MPa$. Eq.~\eqref{eq:zy} puts both reversed limits of hard steel on
the threshold. With $\sigma_x=\sigma_m+\sigma_a\sin\omega t$,
$\tau_{xy}=\tau_m+\tau_a\sin(\omega t-\delta)$:
\begin{center}\small
\begin{tabular}{@{}lccc@{}}
\toprule
Row & load (MPa) & $I$ (\%) & paper\\\midrule
Table~3 row 1, hard steel & $\sigma_a=138.1$, $\tau_a=167.1$, $\delta=0^\circ$ & $-0.09$ & 0.0\\
Table~3 row 4, hard steel & $\sigma_a=150.2$, $\tau_a=181.7$, $\delta=90^\circ$ & $+8.68$ & 8.5\\
Table~4 row 5, 42CrMo4 & $\sigma_a=266$, $\tau_a=128$, $\tau_m=128$, $\delta=0^\circ$ & $-9.85$ & $-9.9$\\
\bottomrule
\end{tabular}
\end{center}

\block{Verification} Tables~3--8, 66 fatigue limits on hard steel, 42CrMo4, 34Cr4, 30NCD16, mild steel
and LY12CZ: the longest-chord amplitude reproduces 50 of 66 rows within 0.3 points (18 within 0.06);
\texttt{rss} reproduces only 6 within 0.06. Eq.~(7) as printed reproduces the 42CrMo4 mean-stress
rows. The maximum-$C_a$ rule is ill-conditioned for this criterion: $C_a$ is stationary at its maximum
while $N^*$ changes to first order, so a $1^\circ$ plane grid moves in-phase indices by up to 0.3
points, and over planes within $10^{-4}$ of the maximum $C_a$ the index spans several points (tens of
points at $90^\circ$); hence the 0.3-point tolerance. Strict xfails: $90^\circ$ rows whose maximum
$C_a$ is shared by a continuum of planes (3-8, 5-3, 5-11, 6-6, 7-6, 8-14, 8-15; the printed value lies
in the tied range), and mean-stress rows of 34Cr4 and 30NCD16 (5-7, 5-8, 5-10, 6-3, 6-4, 6-7 to 6-10),
not reproduced on any near-maximum plane with Eq.~(7) and the $\sigma_u$ of Table~2 (unresolved). Also
checked: Eq.~(6) at both limits, the Goodman correction, the scope and missing-$\sigma_u$ behaviour.
